\section{Recipe decoding and fusion example}
\label{app:recipe-fusion}

We illustrate recipe fusion using two short recipe documents—Italian (cacio e pepe) and Japanese (kake udon), each encoded independently by the Mistral-7B \doctolora hypernetwork. 
We decoded the recipes at both endpoints and their interpolation (full-rank mode, \secref{doc2lora-embedding}), by prompting: ``Write out this recipe: give the dish a name, list the ingredients, and describe the preparation steps.'' 
For example, with only the Italian adapter active, the model answers ``The cheese used in this recipe is Pecorino Romano.'' when asked about the cheese, even without seeing the recipe text. 
Mistral-7B was used (rather than Qwen3-4B) for its more natural recipe outputs, but this does not affect the results.
All outputs below are shown verbatim.

\textbf{Source document $A$ (Italian).}

\begin{snugshade}
\begin{Verbatim}[breaklines=true,breakanywhere=true,fontsize=\footnotesize]
Cacio e Pepe (Roman pasta with cheese and black pepper)

A classic recipe of the Roman tradition with only three ingredients. Serves 2.

Ingredients:
- 200 g spaghetti (or tonnarelli)
- 100 g Pecorino Romano cheese, finely grated
- 2 teaspoons whole black peppercorns
- coarse salt for the pasta water

Preparation:
1. Toast the black peppercorns in a dry skillet over medium heat until
   fragrant, then crush them coarsely in a mortar.
2. Bring a pot of lightly salted water to a boil and cook the spaghetti until
   one minute short of al dente. Use less water than usual so the cooking
   water becomes rich in starch.
3. While the pasta cooks, mix the grated Pecorino Romano with a ladle of the
   hot, starchy pasta water, stirring vigorously until it forms a thick,
   smooth cream with no lumps.
4. Transfer the spaghetti to the skillet with the crushed pepper, add a splash
   of pasta water, and finish cooking for one minute, tossing constantly.
5. Off the heat, add the Pecorino cream and toss vigorously, loosening with
   more pasta water as needed, until every strand is coated in a glossy sauce.
6. Serve immediately with extra grated Pecorino and a final grind of pepper.
\end{Verbatim}
\end{snugshade}

\textbf{Source document $B$ (Japanese).}

\begin{snugshade}
\begin{Verbatim}[breaklines=true,breakanywhere=true,fontsize=\footnotesize]
Kake Udon (Japanese udon noodles in hot dashi broth)

A simple and warming Japanese noodle soup built on dashi stock. Serves 2.

Ingredients:
- 2 portions fresh or frozen udon noodles
- 600 ml dashi stock (kombu and katsuobushi)
- 3 tablespoons soy sauce
- 2 tablespoons mirin
- 1 teaspoon sugar
- 2 spring onions, thinly sliced
- shichimi togarashi (Japanese seven-spice), to taste
- optional toppings: kamaboko fish cake, tempura flakes (tenkasu), a soft egg

Preparation:
1. Make the dashi: soak a piece of kombu in cold water for 30 minutes, bring
   it almost to a boil and remove the kombu, then add a handful of
   katsuobushi (dried bonito flakes), turn off the heat, let it steep for
   2 minutes, and strain.
2. Season the dashi with the soy sauce, mirin, and sugar, and keep the broth
   hot over low heat without boiling.
3. Cook the udon noodles in a separate pot of unsalted boiling water according
   to the package directions, then drain and rinse briefly to remove surface
   starch.
4. Divide the noodles between two deep bowls and ladle the hot seasoned broth
   over them.
5. Top with sliced spring onions and any optional toppings, and finish with a
   pinch of shichimi togarashi. Serve immediately, very hot.
\end{Verbatim}
\end{snugshade}

\textbf{Decoded midpoint, $\alpha=0.5$.}

\begin{snugshade}
\begin{Verbatim}[breaklines=true,breakanywhere=true,fontsize=\footnotesize]
Recipe Name: Kake Udon with Creamy Parmesan Sauce

Ingredients:
- 200 g fresh udon noodles
- 2 tablespoons unsalted butter
- 2 tablespoons mirin
- 2 tablespoons soy sauce
- 1 teaspoon wasabi paste
- 1/2 cup grated Parmesan cheese
- 1/2 cup heavy cream
- 1 spring onion, thinly sliced
- Optional: a pinch of freshly ground black pepper

Preparation Steps:
1. Cook the fresh udon noodles according to the package instructions. Drain
   and rinse with cold water.
2. In a deep skillet or pan, melt the unsalted butter over medium heat.
3. Add the mirin and cook fsr 1 minute.
4. Stir in the soy sauce and wasabi paste.
5. Gradually add the grated Parmesan cheese, stirring constantly to prevent
   lumps from forming.
6. Pour in the heavy cream and bring to a simmer.
7. Cook for 2 minutes, stirring occasionally, until the sauce has thickened.
8. Add the cooked udon noodles to the skillet and toss to coat them evenly
   with the sauce.
9. Remove from heat and top with the sliced spring onion.
10. Optional: finish with a pinch of freshly ground black pepper.
11. Serve immediately.
\end{Verbatim}
\end{snugshade}

The midpoint describes a dish that appears in neither source. Udon noodles coated in a melted-cheese cream sauce, the cacio e pepe technique, seasoned with the mirin and soy of the Japanese broth and finished with the spring onion of one source and the black pepper of the other. The endpoints are faithful. $\alpha=0$ decodes to cacio e pepe and $\alpha=1$ to kake udon, each with minor paraphrases. Off-midpoint interpolants snap to the nearer source ($\alpha=0.25$ decodes to cacio e pepe; $\alpha=0.75$ to kake udon with butter entering the broth).
We note that the blend is imperfect in an instructive way. Namely, the midpoint substitutes Parmesan for Pecorino Romano and introduces heavy cream found in neither source, so the fusion is a plausible blend of the two cuisines rather than a constraint-satisfying merge of the two ingredient lists.
